\begin{figure}[htbp]
  \centering
  \begin{tikzpicture}[x=17mm,y=-17mm,alt={The upper-left two by two room contains cell zero at top left, one at top right, eight at bottom left and nine at bottom right. All other maze cells are walls. Orthogonal links form a cycle. Two routes from zero to nine have two edges: through one or through eight. Arrows highlight the first route; a parent walk traverses it backward from nine to one to zero.}]
    \draw[thick] (0,0) rectangle (2,2);
    \draw (1,0)--(1,2) (0,1)--(2,1);
    \node at (.5,.5) {0}; \node at (1.5,.5) {1};
    \node at (.5,1.5) {8}; \node at (1.5,1.5) {9};
    \draw[->,very thick] (.72,.5)--(1.28,.5);
    \draw[->,very thick] (1.5,.72)--(1.5,1.28);
    \node[anchor=west,align=left,text width=62mm] at (2.4,.55) {Two equal-length routes:\\$0\to1\to9$\\$0\to8\to9$};
    \node[anchor=west,align=left,text width=62mm] at (2.4,1.5) {For the highlighted route,\\parents lead backward:\\$9\to1\to0$.};
  \end{tikzpicture}
  \caption{Cell IDs use the full maze's eight-cell row stride, even in this small open room. A route chooses edges; a parent entry records the preceding cell's ID.}
  \label{fig:revision-maze-room}
\end{figure}
